\begin{abstract}
Automatic interpretation of the 12-lead electrocardiogram (ECG) can speed up the detection of
myocardial infarction, ischaemic changes, conduction disturbances and hypertrophy, but the deep
networks that perform best are trained on centrally pooled data, whereas hospitals cannot share
patient records. Federated learning removes this constraint, yet common methods average model
parameters over many rounds and therefore force every institution to train the same architecture,
and few ECG studies explain their predictions or state when a finding should be referred. This paper
proposes CardioSentinel, a one-shot federated model in which four clients train structurally different
networks (a 1D CNN, a 1D ResNet, a BiLSTM and a CNN-LSTM) on disjoint, patient-separated data, calibrate
them with a temperature, and share only class probabilities with a server that combines them with
weights on the probability simplex. Equal, learned static and per-patient dynamic weights are selected
by patient-grouped cross-validation with a parsimony margin. Experiments use 21,388 records of PTB-XL
for multi-label classification of five superclasses, with 80\% of the data for training, 10\% for
validation and an untouched 10\% for testing, over three seeds. The federated model achieves a macro
AUROC of 0.905, a macro F1 of 70.7\% and a label-wise accuracy of 86.9\%, improves on the best single
client in every seed (+0.009 to +0.012 AUROC, paired bootstrap), and comes within 0.004 AUROC of a
centrally trained model and within 0.002 of FedAvg, which needs 20 rounds and a shared architecture. Integrated Gradients localise the evidence in leads and
waveform segments, and per-finding triage decides 53.7\% of findings automatically at 0.937 precision.
Most of the benefit of pooling ECG data is thus obtained without exchanging data, parameters or
architectures.
\end{abstract}

\begin{IEEEkeywords}
Electrocardiogram, federated learning, ensemble learning, multi-label classification, PTB-XL,
explainable artificial intelligence, calibration.
\end{IEEEkeywords}
